\section{Proofs}
\label{app:proofs}

\begin{proof}[Proof of Proposition~\ref{prop:keep}]
From \eqref{eq:payoff}, $\Pi(K) = (1-b)\sum_i v_i + b \sum_{i \in K} v_i$. The first term does not depend on $K$, so the manager maximizes $\sum_{i \in K} v_i$ subject to $|K| \leq A$. Since the $v_i$ are ordered, the maximizer is the set of the $\min\{n, \lfloor A \rfloor\}$ largest $v_i$.
\end{proof}

\begin{proof}[Proof of Proposition~\ref{prop:cutoff}]
Let $k = |K|$, let $R_k = \sum_{i > k} v_i$ be the importance of the delegated tasks, and let $a = A - k$. Since the manager keeps the most important tasks (by the argument of Proposition~\ref{prop:keep}), her payoff is $\Pi(k) = \sum_i v_i - b e^{-(A-k)} R_k$. Hence
\[
\Pi(k+1) - \Pi(k) = b e^{-A} e^{k} \left[ R_k - e R_{k+1} \right] = b e^{-A} e^{k} \left[ v_{k+1} - (e-1) R_{k+1} \right],
\]
whose sign does not depend on $A$. Let $\hat{k}$ be the smallest $k$ with $v_{k+1} \leq (e-1) R_{k+1}$ (or $\hat{k} = n$ if there is none). Because $R_{k+1}$ decreases in $k$ and $v_{k+1}$ decreases in $k$, the bracket can change sign, so $\hat{k}$ is defined as the first $k$ at which keeping one more task stops paying. The constraint $k \leq A$ then gives $k^* = \min\{\lfloor A \rfloor, \hat{k}\}$, and the number of delegated tasks $n - k^*$ is weakly decreasing in $A$ and constant for $A \geq \hat{k}$.
\end{proof}
